\section{Haplotype alleles}
\label{sec:haplotypes}

The statistics in section~\ref{sec:intro} count alleles at a locus. With
phased data, they can instead count haplotypes across a short run of loci.
\flag{--hap-snps} and \flag{--hap-bp} collapse consecutive SNPs into blocks,
and the alleles of a block are the haplotypes carried by its gene copies.

\begin{Verbatim}[frame=single,fontsize=\small]
adze --data cohort.vcf.gz --samples pops.tsv --hap-snps 5 --out-prefix hap
adze --data cohort.vcf.gz --samples pops.tsv --hap-bp 10000 --out-prefix hap
\end{Verbatim}

\noindent The estimators do not change. $\hat{\alpha}_g^{(j)}$ and
$\hat{\pi}_g^{(j)}$ are sums over the alleles at a locus, and a block is
treated as a locus whose alleles are haplotypes. The rarefaction, the
averaging across loci and the output are the same as described in the rest of
this manual. However, because an allele is now a haplotype, there are some
additional requirements and caveats, which are described below.

\subsection{What a haplotype allele is here}

For each gene copy, \prog\ reads its allele at every SNP in the block, in order
along the chromosome, and treats this sequence as one allele. Two gene copies
carry the same haplotype allele if they match at every SNP in the block. The
blocks do not overlap, so each SNP belongs to exactly one block and no gene
copy is counted twice.

\flag{--hap-snps} $N$ makes each block $N$ consecutive SNPs, so every block has
the same number of SNPs, and its width in basepairs varies with the density of
SNPs. \flag{--hap-bp} $N$ makes each block a region of $N$ basepairs, starting
at position~1 and moving by $N$. The same region therefore gives the same
blocks in every run, so two data sets can be compared block by block, and the
number of SNPs in a block varies instead. A block never spans two
chromosomes, a region that contains no SNP gives no block, and
\flag{--min-hap-snps} drops blocks that contain too few SNPs.

\subsection{Phase is required}

A haplotype is a run of alleles along one chromosome, so the input must say
which chromosome copy each allele is on. \prog\ therefore requires phased
genotypes. Every \texttt{GT} with more than one allele must use \texttt{|}, and
a record that uses \texttt{/} stops the run with an error that gives the record
and the sample. Haploid gene copies are phased by definition. Note that \prog\
does not phase the data or check the quality of the phasing. A haplotype
assembled from an unphased genotype would be an arbitrary pairing of alleles,
which is why an unphased genotype is an error rather than a warning.

For the same reason, the \structureprog\ format is not supported in this mode,
since it does not say which chromosome copy an allele is on. Giving
\flag{--hap-snps} or \flag{--hap-bp} with \structureprog\ input is refused as a
usage error.

The records must also be sorted by chromosome and then by position. A block is
a run of consecutive records, so if the records were out of order a block would
join SNPs that are not next to each other. An ordinary run only needs sorted
input for sliding windows, but a haplotype run always needs it. The run stops
at the first record whose position is smaller than the record before it, or
that returns to a chromosome that has already ended (\texttt{bcftools sort}
fixes either problem).

\subsection{One uncalled genotype makes the block uncalled}

A gene copy only carries a haplotype allele in a block if it is called at every
SNP in the block. If it is uncalled at any SNP in the block, it does not
contribute to that block. It is not given a partial haplotype, since this
would be an allele that no gene copy carries, and it would increase every
count that included it. A gene copy whose \texttt{GT} at some record has fewer
alleles than the sample's ploidy is also counted as uncalled there.

As a result, missing data accumulate along a block. If each SNP is called
independently with probability $1-m$, a gene copy is callable across a block of
$N$ SNPs with probability about $(1-m)^N$. For example, with 2\% missing data,
a block of 50 SNPs keeps about a third of its gene copies. Because $N_j$ enters
the rarefaction directly, this reduces the largest $g$ the data can support.
\flag{--tolerance} applies to blocks in the same way as to loci, so a block
where some grouping has lost more than the allowed fraction of its gene copies
is dropped. The \texttt{\_blocks} file (section~\ref{sec:blocksfile}) gives
the number of callable gene copies in each grouping for every block.

\subsection{Richness saturates as blocks lengthen}

When a block is long enough, every gene copy carries a different haplotype.
If a block has as many distinct haplotypes as callable gene copies, no
haplotype is repeated, every subsample of $g$ gene copies contains $g$ distinct
haplotypes, and $\hat{\alpha}_g^{(j)} = g$ regardless of the data. In this case
the statistic no longer measures anything about the populations. \prog\ counts
these blocks and reports the number at the end of the run, and the
\texttt{HAPLOTYPES} column of the \texttt{\_blocks} file shows how close each
block is to this limit.

Private allelic richness has the same limit, because a haplotype carried by
only one gene copy is private to that gene copy's grouping. For example, in
simulated biallelic data in which both alleles at every SNP have a frequency
between 0.2 and 0.8 in every grouping, so that almost no single-SNP allele is
private, the mean private allelic richness at $g = 2$ was 0.0025 for blocks of 1 SNP, 0.11 for 2 SNPs, 1.42 for 5 SNPs, 1.98
for 10 SNPs and 2.0 for 20 or 50 SNPs. We therefore recommend using the
shortest blocks that answer the question, and comparing results across several
block lengths.

\subsection{Naming and output}
\label{sec:blocksfile}

A block with more than one SNP is named \texttt{chromosome:start-end}. A block
with exactly one SNP takes the name of that record, so \flag{--hap-snps 1}
gives exactly the same output as an ordinary run. All of the output files
described in section~\ref{sec:output} are written, with blocks in place of
loci. This includes the \texttt{\_fulldata} files, and the window files in a
windowed run.

One additional file, \texttt{\_blocks}, describes the blocks themselves:

\begin{Verbatim}[frame=single,fontsize=\footnotesize]
BLOCK           CHROM START END   SNPS HAPLOTYPES KEPT POP0 POP1 POP2 POP3
chr1:1000-4000  chr1  1000  4000  4    14         no   9    5    8    8
chr1:5000-8000  chr1  5000  8000  4    11         yes  12   8    7    8
\end{Verbatim}

\noindent There is one row for each block, in the order of the input file. Each
row gives the number of callable gene copies in each grouping and whether the
block was kept by the filter. Two runs can give the same mean richness with
blocks that have very different numbers of haplotypes, and this file shows the
difference.

\subsection{Haplotype blocks in sliding windows}

Haplotype blocks can be used with sliding windows. Each window averages over
the blocks that it contains, in the same way that an ordinary windowed run
averages over loci (section~\ref{sec:windows}), and \flag{--window-loci}
counts blocks. A block is placed in a window by its first SNP. This has two
consequences.

With \flag{--window-loci}, the reported interval of a window runs from the
first SNP of its first block to the last SNP of its last block, so it covers
all of the SNPs in the window.

With \flag{--window-bp}, a block belongs to the window that contains its first
SNP, and it can extend past the end of that window. Its haplotypes then
include SNPs outside the reported interval of the window. \prog\ counts these
blocks and reports the number. Blocks always fall inside windows when
\flag{--hap-bp} $h$ divides both \flag{--window-bp} and \flag{--step-bp},
since both are counted from position~1. For example, the following run has no
blocks that extend past their window:

\begin{Verbatim}[frame=single,fontsize=\small]
adze --data cohort.vcf.gz --samples pops.tsv --hap-bp 5000 \
     --window-bp 100000 --step-bp 25000 --out-prefix scan
\end{Verbatim}

\noindent Blocks from \flag{--hap-snps} do not have a fixed width, so they
cannot be made to fall inside windows in this way. With typical SNP densities,
a block is short compared with a window, and only the blocks at the edge of a
window are affected.

\flag{--dry-run} reports the blocks that a given block size produces, the
number of SNPs they contain and their widths. It uses only the first pass over
the data, so it is a fast way to choose a block size before running a full
genome scan.
